\rtmtight
\begin{tblr}{width=\textwidth,colspec={|Q[c,m,2.1cm]|X[l,m]|},hlines,vlines,hspan=minimal,rowsep=1.5pt,colsep=3pt,row{1}={bg=headgray,font=\bfseries}}
\SetCell{c,m}\hfil\logo\hfil & \textbf{POLITEKNIK NEGERI MALANG}\par\textbf{JURUSAN TEKNOLOGI INFORMASI}\par\textbf{PROGRAM STUDI: D-IV SISTEM INFORMASI BISNIS} \\
\SetCell[c=2]{l} \textbf{RENCANA TUGAS MAHASISWA (RTM) DAN RUBRIK PENILAIAN} & \\
Nama Mata Kuliah & Pemrograman Web Lanjut \\
Kode Mata Kuliah & SIB245007 \quad \textbf{sks / Jam perminggu:} 3 SKS/6 jam \quad \textbf{Semester:} 4 \\
Dosen Pengampu & \begin{enumerate}\item Dian Hanifudin Subhi, S.Kom., M.Kom.\item Ade Ismail, S.Kom., M.TI.\end{enumerate} \\
\end{tblr}

\assessmentblock{Tugas Praktik: Framework, Routing, \& Basis Data}{Hands-on Practice}{SCPMK0202-04001, SCPMK0202-04002}{Mahasiswa menyelesaikan tugas praktik mingguan terkait setup framework, routing, controller, desain basis data, migrasi skema, ORM, dan server-side templating; setiap tugas merupakan increment yang membangun aplikasi studi kasus POS/kasir UMKM secara bertahap.}{Menganalisis kebutuhan, mengerjakan artefak praktik, menguji hasil, mendokumentasikan evidence, dan mempresentasikan keputusan bila diminta.}{Laporan, artefak digital, repository/prototype, evidence pengujian, dan/atau bahan presentasi.}{Kesiapan lingkungan \& konfigurasi framework & 4 & Framework terpasang, konfigurasi basis data, struktur proyek, dan version control siap digunakan secara penuh. & Framework berjalan tetapi konfigurasi atau struktur awal belum lengkap. & Framework tidak berjalan atau konfigurasi dasar salah. \\ \\
Routing, controller, basis data, migrasi skema, \& ORM & 3 & Route, controller, migrasi skema, model, relasi ORM, dan accessor/mutator dasar sesuai kebutuhan modul serta konsisten. & Komponen utama tersedia tetapi relasi ORM, migrasi, atau controller belum sepenuhnya konsisten. & Routing/model/migrasi/ORM tidak terhubung atau tidak sesuai fitur. \\ \\
Templating, validasi, \& evidence praktik & 3 & Server-side template, formulir, validasi, feedback error, dan screenshot/commit praktik terdokumentasi jelas. & Template dan formulir berjalan tetapi validasi atau evidence praktik masih terbatas. & Template tidak mendukung alur fitur atau tidak ada bukti praktik. \\}

\assessmentblock{UTS: Evaluasi Milestone Aplikasi Studi Kasus}{UTS praktik}{SCPMK0202-04001, SCPMK0202-04002, SCPMK0202-04003}{Mahasiswa mendemonstrasikan milestone aplikasi studi kasus POS/kasir UMKM (hasil increment minggu 1--7) yang mencakup CRUD lengkap, autentikasi, otorisasi, dan arsitektur MVC, kemudian mempertahankan keputusan teknis secara lisan (viva). Evaluasi bersifat individu, bukan proyek kelompok PBL.}{Menganalisis kebutuhan, mengerjakan artefak proyek, menguji hasil, mendokumentasikan evidence, dan mempertahankan keputusan teknis secara lisan.}{Aplikasi web berjalan, repository, evidence pengujian, dan/atau bahan viva.}{Kelengkapan CRUD \& validasi & 2 & Create, read, update, delete, validasi input, flash message, dan penanganan error berjalan lengkap. & CRUD utama berjalan tetapi validasi atau penanganan error belum lengkap. & CRUD tidak berjalan utuh atau data tidak tersimpan benar. \\ \\
Kualitas arsitektur MVC \& basis data & 2 & Route, controller, model, migrasi, seeder, dan view dipisahkan rapi mengikuti konvensi framework yang digunakan. & Struktur cukup mengikuti arsitektur MVC tetapi masih ada logika bercampur atau konvensi tidak konsisten. & Kode tidak mengikuti arsitektur MVC atau sulit dipelihara. \\ \\
Pengujian praktik \& defense & 1 & Skenario CRUD normal/invalid diuji, data diverifikasi di basis data, dan keputusan teknis dapat dijelaskan secara lisan. & Pengujian dilakukan terbatas atau defense belum menjelaskan semua keputusan. & Tidak ada pengujian bermakna atau tidak mampu menjelaskan modul. \\}

\assessmentblock{Tugas Praktik: Autentikasi, API, \& Pengolahan Data}{Hands-on Practice}{SCPMK0202-04003, SCPMK0202-04004}{Mahasiswa menyelesaikan tugas praktik terkait autentikasi, otorisasi berbasis peran, desain API, token-based authentication, konsumsi API dari sisi klien, serta pengolahan dan ekspor data; setiap tugas merupakan increment lanjutan yang menyelesaikan aplikasi studi kasus POS/kasir UMKM.}{Menganalisis kebutuhan, mengerjakan artefak praktik, menguji hasil, mendokumentasikan evidence, dan mempresentasikan keputusan bila diminta.}{Laporan, artefak digital, repository/prototype, evidence pengujian, dan/atau bahan presentasi.}{Autentikasi \& otorisasi berbasis peran & 2 & Login, logout, manajemen peran/hak akses, proteksi route, dan akses data diterapkan sesuai skenario. & Autentikasi berjalan tetapi manajemen peran, proteksi route, atau akses data masih memiliki celah. & Autentikasi/otorisasi tidak berjalan atau route penting tidak terlindungi. \\ \\
Desain API \& pengolahan data & 1 & Endpoint API, request/response, validasi, penanganan error, import/export data, dan token-based auth bekerja konsisten. & Fitur API/pengolahan data berjalan sebagian tetapi validasi atau penanganan error belum lengkap. & API gagal atau fitur import/export tidak sesuai kontrak data. \\ \\
Konsumsi API dari sisi klien & 1 & Contoh klien (consumer) yang mengonsumsi API bekerja dengan benar: berhasil login, menyimpan/mengirim token, memanggil endpoint terproteksi, dan menangani respons/galat dengan tepat. & Klien berhasil memanggil sebagian endpoint tetapi penanganan token atau galat belum konsisten. & Klien tidak berhasil mengonsumsi API atau tidak ada contoh klien sama sekali. \\ \\
Dokumentasi \& evidence pengujian & 1 & Dokumentasi endpoint (OpenAPI/Swagger atau setara), contoh request/response, dan evidence eksekusi tersedia jelas. & Dokumentasi tersedia tetapi cakupan dan evidence masih terbatas. & Tidak ada dokumentasi memadai atau evidence tidak membuktikan fitur. \\}

\assessmentblock{PBL: Persiapan Tim \& Eksplorasi Domain}{Project Based Learning}{SCPMK0202-04004}{Kelompok PBL memulai proyek aplikasi web pada minggu 1--3: membentuk tim, membagi peran awal, dan menjajaki domain bisnis pilihan tim (mis. inventori, koperasi, reservasi, laundry) sebagai persiapan menuju pengajuan proposal pada minggu 4.}{Membentuk tim, membagi peran awal, menjajaki dan membandingkan kandidat domain bisnis, serta mengidentifikasi masalah pengguna pada domain terpilih.}{Catatan pembentukan tim \& peran, hasil eksplorasi domain, dan evidence diskusi/rapat kelompok setiap minggu.}{Pembentukan tim \& pembagian peran awal & 2 & Tim terbentuk lengkap dengan peran/tanggung jawab awal yang jelas serta rencana kerja kelompok tercatat pada minggu 1. & Tim terbentuk tetapi peran atau rencana kerja belum sepenuhnya jelas. & Tim belum terbentuk atau tidak ada pembagian peran. \\ \\
Eksplorasi domain bisnis \& identifikasi masalah & 2 & Minimal dua kandidat domain bisnis dieksplorasi dengan identifikasi masalah pengguna yang jelas pada minggu 2. & Eksplorasi domain dilakukan tetapi masalah pengguna belum teridentifikasi jelas. & Tidak ada evidence eksplorasi domain bisnis. \\ \\
Pemilihan domain bisnis \& kesiapan menuju proposal & 2 & Domain bisnis final dipilih dengan justifikasi jelas dan tim siap menyusun proposal pada minggu 3. & Domain dipilih tetapi justifikasi masih lemah atau kesiapan proposal belum matang. & Domain bisnis belum ditentukan menjelang minggu 4. \\}

\assessmentblock{PBL: Proposal Proyek}{Project Based Learning}{SCPMK0202-04004}{Kelompok PBL menyusun dan mengajukan proposal proyek pada minggu 4, memuat domain bisnis terpilih, kebutuhan pengguna, model data awal, dan daftar fitur awal sebagai dasar pengembangan aplikasi sepanjang semester.}{Menganalisis kebutuhan pengguna, merancang model data awal, menyusun daftar fitur awal, mendokumentasikan proposal, dan mempresentasikannya.}{Dokumen proposal proyek (minggu 4) dan/atau bahan presentasi proposal.}{Kejelasan domain bisnis \& masalah yang diselesaikan & 4 & Domain bisnis dan masalah pengguna dijelaskan spesifik, relevan, dan didukung justifikasi yang kuat. & Domain bisnis dijelaskan tetapi masalah pengguna kurang spesifik atau justifikasi lemah. & Domain bisnis atau masalah pengguna tidak jelas. \\ \\
Kelengkapan kebutuhan pengguna \& model data awal & 4 & Kebutuhan pengguna terpetakan lengkap dan model data awal (entitas \& relasi) konsisten dengan kebutuhan tersebut. & Sebagian kebutuhan pengguna atau model data awal belum terpetakan jelas. & Kebutuhan pengguna tidak terpetakan atau model data tidak ada. \\ \\
Realistis \& terukur daftar fitur awal & 4 & Daftar fitur awal realistis untuk satu semester, terukur, dan diprioritaskan sesuai kebutuhan pengguna. & Daftar fitur awal ada tetapi prioritas atau keterukurannya belum jelas. & Daftar fitur awal tidak realistis atau tidak ada. \\ \\
Kualitas dokumentasi \& kesiapan presentasi proposal & 3 & Dokumen proposal tersusun rapi, lengkap, dan dipresentasikan dengan jelas pada minggu 4. & Dokumen proposal ada tetapi kurang rapi atau presentasi kurang meyakinkan. & Dokumen proposal tidak diserahkan atau tidak dipresentasikan. \\}

\assessmentblock{PBL: Pengembangan \& Finalisasi Proyek}{Project Based Learning}{SCPMK0202-04004}{Kelompok PBL mengembangkan proyek aplikasi web berbasis framework secara terintegrasi berdasarkan proposal minggu 4, terinspirasi dari aplikasi studi kasus POS/kasir UMKM. Pengembangan berlanjut bertahap di sela-sela tugas individu (minggu 5--7, 9--10), didetailkan menjadi rencana teknis (minggu 11), dikerjakan intensif pada minggu proyek khusus (minggu 12--14), lalu difinalisasi dan direhearsal (minggu 15--16) sebelum dipresentasikan pada UAS.}{Merencanakan arsitektur, mengerjakan artefak proyek secara bertahap, menguji hasil, mendokumentasikan evidence, dan mempresentasikan produk.}{Laporan proyek, artefak digital, repository, video demo, dan/atau bahan presentasi.}{Kontribusi progress mingguan (minggu 5--7, 9--10) & 12.5 & Progress mingguan proyek terdokumentasi konsisten dengan fitur bertambah sesuai proposal pada setiap minggu 5, 6, 7, 9, dan 10. & Progress mingguan ada tetapi tidak konsisten atau beberapa minggu tanpa evidence. & Tidak ada evidence progress mingguan proyek. \\ \\
Perencanaan teknis lanjutan & 3 & Rencana teknis (arsitektur aplikasi terperinci \& pembagian peran tim) mendetailkan proposal secara konsisten pada minggu 11. & Rencana teknis ada tetapi belum sepenuhnya konsisten dengan proposal. & Tidak ada rencana teknis lanjutan pada minggu 11. \\ \\
Pengembangan fitur utama \& integrasi (minggu 12--14) & 10 & CRUD, autentikasi, API/import-export, validasi, dan relasi basis data terintegrasi secara stabil pada minggu proyek khusus. & Fitur utama terintegrasi tetapi masih ada kekurangan validasi, relasi, atau penanganan error. & Fitur utama tidak terintegrasi atau aplikasi sering gagal. \\ \\
Finalisasi \& rehearsal (minggu 15--16) & 8.5 & Pengujian, data demo/seeder, deployment atau demo lokal, repository, dan rehearsal/mock defense minggu 16 tersedia lengkap. & Evidence finalisasi tersedia tetapi pengujian/deployment atau rehearsal belum kuat. & Tidak ada evidence memadai atau proyek tidak dapat didemonstrasikan. \\}

\assessmentblock{UAS: Evaluasi Proyek Aplikasi Web Final}{UAS praktik dan defense}{SCPMK0202-04004}{Mahasiswa mendemonstrasikan proyek web final, mempertahankan keputusan teknis secara lisan, dan mengevaluasi ketercapaian kebutuhan pengguna.}{Menganalisis ketercapaian kebutuhan, mendemonstrasikan aplikasi, mempertahankan keputusan teknis, dan memberikan rekomendasi peningkatan.}{Aplikasi web final, repository, dokumentasi teknis, dan/atau bahan presentasi/defense.}{Evaluasi kelengkapan fitur akhir & 10 & Fitur aplikasi memenuhi scope, peran, alur kebutuhan pengguna, validasi, dan pengelolaan data secara menyeluruh. & Fitur utama memenuhi scope tetapi beberapa validasi atau peran masih belum lengkap. & Fitur akhir tidak memenuhi scope atau banyak alur gagal. \\ \\
Kualitas teknis \& maintainability & 9 & Kode mengikuti konvensi framework, struktur rapi, query efisien, dan konfigurasi siap dijalankan ulang. & Kode cukup rapi tetapi masih ada duplikasi, query kurang efisien, atau konfigurasi kurang jelas. & Kode sulit dipelihara atau tidak dapat dijalankan ulang. \\ \\
Defense, evidence, \& rekomendasi & 6 & Defense menjelaskan arsitektur, pengujian, deployment, risiko, dan rekomendasi peningkatan aplikasi secara jelas. & Defense cukup jelas tetapi evidence pengujian/deployment atau rekomendasi masih terbatas. & Tidak mampu mempertahankan keputusan teknis atau evidence tidak memadai. \\}

\begin{tblr}{width=\textwidth,colspec={|Q[l,m,3.2cm]|X[l,m]|},hlines,vlines,hspan=minimal,row{1-3}={bg=headgray},rowsep=1.5pt,colsep=3pt}
Jadwal Pelaksanaan: & Minggu 1--17 sesuai RPS. Setiap evaluasi memiliki RTM dan rubrik multi-kriteria. \\
Lain-Lain Yang Diperlukan: & LMS, repositori/prototype/proyek, perangkat lunak sesuai mata kuliah, dan perangkat presentasi. \\
Pustaka: & Mengacu pustaka utama dan pendukung pada bagian RPS. \\
\end{tblr}
